\section{Method}

\subsection{Overview}

Our goal is to reconstruct a high-fidelity object-centric 3D representation from a monocular video. Given an input video sequence $V={I_i}_{i=1}^{N}$, we first estimate camera parameters $C={K_i,T_i}_{i=1}^{N}$ using COLMAP and obtain foreground masks $M={M_i}_{i=1}^{N}$ with SAM2.

Building upon these semantic priors, we propose a reconstruction framework that combines mask-guided optimization with geometry-aware refinement. During training, foreground masks are incorporated into the 3D Gaussian Splatting (3DGS) optimization process to suppress background fitting. After reconstruction, we apply a multi-metric pruning strategy based on multi-view consistency and local geometric properties to remove floating artifacts and refine object boundaries. The overall pipeline is illustrated in Figure~\ref{fig:schematic}.

\begin{figure}[t]
    \centering
    \includegraphics[width=1\linewidth]{pics/MOF3R_overview.png}
    \caption{MOF3R: Overall pipeline schematic}
    \label{fig:schematic}
\end{figure}

\subsection{Mask-Guided Optimization}

A major challenge of object-centric reconstruction from unconstrained videos is that standard 3DGS optimizes all image pixels equally, causing Gaussian primitives to fit both the target object and surrounding background regions.

To focus optimization on the object of interest, we incorporate foreground masks generated by SAM2 into the training objective. Given a ground-truth image $I$, rendered image $\hat I$, and binary mask $M$, we define a mask-constrained L1 loss as

$$
\mathcal{L}_{L1}^{mask} = |M\odot(I-\hat I)|_1
$$

where $\odot$ denotes multiplication by element. This formulation restricts photometric supervision to foreground regions and suppresses background fitting during the early stages of optimization.

Directly applying structural similarity loss only inside the mask may introduce instability near object boundaries. Therefore, we construct a composite image

$$
I_{comp} = M\odot\hat I + (1-M)\odot I
$$

where rendered foreground pixels are blended with the original background. The SSIM loss is then computed on the composite image:

$$
\mathcal{L}_{\mathrm{SSIM}}^{mask} = 1-\mathrm{SSIM}(I, I_{comp})
$$

The final training objective is defined as

$$
\mathcal{L} = \lambda_1\mathcal{L}_{L1}^{mask} + \lambda_2\mathcal{L}_{\mathrm{SSIM}}^{mask}
$$

This compound objective encourages Gaussian primitives to concentrate on the target object while preserving global image consistency, resulting in cleaner reconstructions and fewer background artifacts.

\subsection{Geometry-Aware Multi-Metric Pruning}

Although mask-guided optimization effectively suppresses most background interference, the reconstructed model may still contain floating Gaussians and elongated boundary artifacts caused by limited viewpoints and optimization errors. To further improve reconstruction quality, we introduce a geometry-aware multi-metric pruning strategy.

\subsubsection{Multi-view Mask Consistency} 

For each Gaussian center $g_j$, we project its 3D position into all visible camera views and examine whether the projected point lies within the corresponding foreground masks.

Gaussians that consistently fall outside foreground regions across multiple valid views are considered background artifacts and removed. To avoid incorrect pruning caused by occlusions, only visible views participate in the voting process.

This consistency check effectively eliminates isolated floating Gaussians around the reconstructed object.

\subsubsection{Density and Anisotropy Filtering}

Multi-view consistency alone is insufficient to remove all boundary artifacts. Therefore, we further analyze the local geometric structure of each Gaussian using a K-nearest-neighbor (KNN) neighborhood.

For each Gaussian primitive, we estimate its local density and anisotropy level from neighboring Gaussians. A Gaussian is marked as an outlier if it simultaneously exhibits:

\begin{enumerate}
    \item significantly lower local density than its neighborhood;
    \item excessive anisotropy beyond a predefined threshold.
\end{enumerate}

Such primitives typically correspond to stretched structures and trailing artifacts near object boundaries. Removing them leads to a cleaner and more compact reconstruction.

\subsubsection{Adaptive Shrinkage}

Directly removing all boundary candidates may damage fine structures. To preserve delicate geometric details, we introduce an adaptive shrinkage mechanism.

Instead of immediately deleting uncertain boundary Gaussians, we gradually reduce their scale and opacity:

$$
s'=\alpha s,\qquad
o'=\beta o
$$

where $0<\alpha,\beta<1$.

This progressive refinement smooths object boundaries while retaining high-confidence geometric structures, resulting in a cleaner and visually sharper reconstruction.
